\documentclass[10pt,aspectratio=169]{beamer}

% All the boilerplate is in ccaslides.sty
% Note that this also pulls in a custom vogtwidebar.sty
\usepackage{ccaslides}

\author{Ji\v{r}\'i Lebl}

\institute[OSU]{%
Departemento pri Matematiko de Oklahoma {\^S}tata Universitato}

\title{Cultivating Complex Analysis:\\%
Harmonic functions\\%
Isolated singularities (7.3.1)}

\date{}

\begin{document}

\begin{frame}
\titlepage
\end{frame}

\begin{frame}
What can we say about isolated singularities of harmonic functions?

\medskip
\pause

$\log \sabs{z}$ is harmonic with an isolated singularity at $0$.

\medskip
\pause

That's the ``slowest rate'' a harmonic $f$ can blow up at.

\begin{theorem}
Suppose $U \subset \C$ is open, $p \in U$, and $f \colon U \setminus \{ p \}
\to \R$ is harmonic such that
\begin{equation*}
\lim_{z\to p} \frac{f(z)}{\log \sabs{z-p}} = 0 .
\end{equation*}
Then there exists a harmonic $F \colon U \to \R$ such that
$f = F|_{U \setminus \{ p \}}$.
\end{theorem}

\medskip
\pause

\textbf{Proof:}
WLOG, assume (using $f(a z + b)$) that $p = 0$ and $\overline{\D} \subset U$.

\medskip
\pause

Find a continuous $u \colon \overline{\D} \to \R$,
harmonic in $\D$, such that $u|_{\partial \D} = f|_{\partial \D}$
(solve Dirichlet).

\medskip
\pause

If we show that $u = f$ in $\D \setminus \{ 0 \}$, we are done.

\end{frame}

\begin{frame}
We have
$f$ given on $\overline{\D} \setminus \{ 0 \}$,
$\displaystyle \lim_{z\to 0} \frac{f(z)}{\log \sabs{z}} = 0$,
$u$ on $\overline{\D}$, harmonic in $\D$, and $f = u$ on $\partial \D$.

\medskip
\pause

$g = f-u$ is harmonic in $\D \setminus \{ 0 \}$, continuous
on $\overline{\D} \setminus \{ 0 \}$, zero on $\partial \D$, and
\[
\lim_{z \to 0} \frac{g(z)}{-\log \sabs{z}} = 0.
\]
\pause
For any $\epsilon > 0$, there is a $\delta > 0$ so that
\[
-\epsilon (- \log\sabs{z})
\leq
g(z)
\leq
\epsilon (- \log\sabs{z})
\qquad
\text{for all }
z \in \overline{\Delta_\delta(0)} \setminus \{ 0 \}.
\tag{$\dagger$}
\]
(we can assume $\delta < 1$).

\medskip
\pause

The estimate ($\dagger$) holds on $\partial \D$ ($g=0$ there).

\pause
\medskip

$-\log\sabs{z}$ and $g$ are harmonic (outside $0$) \wthus

maximum principle
\wthus
($\dagger$) holds when
$\delta < \sabs{z} < 1$ \wthus ($\dagger$) holds
for all $z \in \D \setminus \{ 0 \}$.

\medskip
\pause

($\dagger$) holds for all $\epsilon > 0$ for all $z \in \D \setminus \{0\}$
\wthus 
$g = 0$ and so $f=u$ on $\D \setminus \{ 0 \}$.
\qed
\end{frame}

\begin{frame}

An isolated singularity of a harmonic $f$ could be wild, e.g., $\Re
e^{1/z}$.  But what if $f \geq 0$?

\medskip
\pause

If $f(z) = \log \sabs{h(z)}$ for a holomorphic $h$
with a zero or pole at the origin,

then near the origin $h(z) = \psi(z) z^n$ where $\psi(0)\not=0$ and $n \in \Z$,
and so
\[
f(z) = 
\log \sabs{\psi(z) z^n} = 
\log \sabs{\psi(z)} + \log \abs{z^n} = 
\log \sabs{\psi(z)} +
n \log \sabs{z} .
\]

So $f$ is a harmonic plus a multiple of $\log \sabs{z}$.

\medskip
\pause

\textit{Remark:}
Near $0$,
$f(z) \geq 0$ ($n < 0$) if $h$ has a pole
and
$f(z) \leq 0$ ($n > 0$) if $h$ has a zero.

\pause

\begin{theorem}[B\^{o}cher]
Suppose $U \subset \C$ is open, $p \in U$, and
$f \colon U \setminus \{ p \} \to \R$ is
harmonic and nonnegative.  Then there exists a harmonic function
$g \colon U \to \R$ and a $C \geq 0$ such that
for all $z \in U \setminus \{ p \}$,
\begin{equation*}
f(z) = g(z) - C \log \sabs{z-p} .
\end{equation*}
\end{theorem}

\end{frame}

\begin{frame}

\textbf{Proof:}
WLOG, $p=0$ and $U = \D$.

\medskip
\pause

$\D \setminus \{ 0 \}$ not simply connected
\wthus
holomorphic $\Phi$ so that $\Re \Phi = f$ not well-defined.

\medskip
\pause

But harmonic conjugates differ by a constant \wthus $\Phi'$ is well-defined.

\medskip
\pause

$\exists$ a holomorphic $\varphi \colon \D \setminus \{ 0 \} \to \C$ such
that $\Phi' = \varphi$ for any such locally defined $\Phi$.

\medskip
\pause

Write the Laurent series:
\quad
$\displaystyle
\varphi(z) = \sum_{n=-\infty}^\infty c_n {z}^n$.

\medskip
\pause

Near some point, for some branch of $\log z$, and a constant $A$ (antidifferentiate):
\[
\Phi(z) =
A+
c_{-1} \log z +  
\sum_{n=-\infty,n \neq -1}^\infty \frac{c_n}{n+1} {z}^{n+1}
\pause
=
c_{-1} \log z +  
\psi(z),
\]
where $\psi$ is holomorphic
on $\D \setminus \{ 0 \}$
and we just assume $A=0$.


\medskip
\pause

$f = \Re \Phi$ and so $\Re (c_{-1} \log z)$ is well-defined on $\D \setminus
\{ 0 \}$
\wthus $c_{-1} \in \R$.

\end{frame}

\begin{frame}
So
\[
f(z) = \Re \Phi(z) =
c_{-1} \log \sabs{z} +  
\Re \psi(z).
\]
\pause
$f(z) \geq 0$ \wthus
$-c_{-1} \log \sabs{z} \leq  \Re \psi(z)$.
\pause
\qquad
Take an integer $k \leq c_{-1}$, then
\[
-k \log \sabs{z} \leq
-c_{-1} \log \sabs{z} \leq  \Re \psi(z) .
\]
\pause
Then
\[
\sabs{z^{-k}} \leq e^{\Re \psi(z)} = \babs{e^{\psi(z)}} 
\pause
\quad
\text{or}
\quad
\sabs{z^{-k} e^{-\psi(z)}} \leq 1 .
\]
\pause
$z^{-k} e^{-\psi(z)}$ has
a removable singularity at the origin.

\medskip
\pause

An exercise from section 5.2.3 says that
$e^{h(z)}$ has removable or essential singularities only.

\pause
\medskip

So $e^{-\psi(z)}$, and so $\psi(z)$, has a removable singularity
\wthus
$\Re \psi(z)$ extends through $0$.

\medskip
\pause

$-c_{-1} \log\sabs{z} \leq \Re \psi(z)$
\wthus
$-c_{-1} \log\sabs{z}$ is bounded from above near $0$
\wthus
$-c_{-1} \geq 0$.

\medskip
\pause

We are done:
\[
f(z) = \Re \psi(z) - (-c_{-1}) \log \sabs{z} . \qed
\]

\end{frame}

\begin{frame}

\textbf{Exercise:}
Prove that the Dirichlet problem is not necessarily solvable in the punctured disc $\D
\setminus \{ 0 \}$.

\medskip

\textit{Remark:} $\partial (\D \setminus \{ 0 \}) = (\partial \D ) \cup \{ 0 \}$.

\medskip
\pause

\textbf{Exercise:}
Prove that given $f$, the function $g$ and the constant $C$ in 
B\^{o}cher's theorem are unique.

\end{frame}

\end{document}
